\documentclass[11pt]{article}
\usepackage[letterpaper,margin=1in]{geometry}
\usepackage{amsmath,amssymb,booktabs}
\setlength{\parskip}{0.5em}\setlength{\parindent}{0pt}
\title{Why temperature only rescales time:\\ the model in dimensionless form}
\author{enceladus\_DSM \quad \texttt{studies/keff\_sintering/temperature\_scaling/}}
\date{2026-10-12}
\begin{document}
\maketitle

\textbf{Result.} For isothermal sintering by vapor transport, the model
reduces to a problem with a single parameter, $\Lambda = L^{*}/\bar R$, once
time is measured in units of $\tau = \bar R^{2}\beta_{\mathrm{sub}}/d_0$. The
saturation vapor density, which changes 31-fold between $-40$ and
$-5\,^{\circ}$C, is entirely inside $\tau$. $\Lambda$ contains the vapor
diffusivity and changes by 20\%. Numbers are from
\texttt{check\_groups.py} (output \texttt{groups.txt}).

\section{The dimensional problem}

The phase-field model is used in its sharp-interface limit, which is the limit
it is built to reproduce (Kaempfer \& Plapp, 2009). The aggregate runs are
isothermal, so temperature is a constant parameter $T$. Three statements
remain.

\emph{Vapor in the pore.} From Eq.~(1c) of the manuscript, away from the
interface,
\begin{equation}
\frac{\partial \rho_v}{\partial t} = D_v \nabla^{2}\rho_v .
\label{eq:diff}
\end{equation}

\emph{Mass balance at the interface.} Let $\mathbf n$ be the unit normal
pointing from the ice into the pore and $v_n$ the velocity of the interface
along $\mathbf n$ ($v_n > 0$ when ice grows). Ice that grows takes its mass
from the vapor arriving by diffusion. Integrating Eq.~(1c) across the
interface, and neglecting $\rho_v$ against $\rho_{\mathrm{ice}}$ (a ratio of
$10^{-6}$),
\begin{equation}
\rho_{\mathrm{ice}}\, v_n = D_v\, \frac{\partial \rho_v}{\partial n} .
\label{eq:stefan}
\end{equation}
The sign: above a growing surface the vapor is depleted, so $\rho_v$ rises
with distance from the ice and $\partial\rho_v/\partial n > 0$.

\emph{Gibbs--Thomson condition at the interface} (manuscript Eq.~6),
\begin{equation}
\frac{\rho_v - \rho_{vs}^{I}}{\rho_{vs}^{I}} = d_0\,\chi + \beta\, v_n ,
\qquad
\beta = \frac{1}{\alpha_c}\,\frac{\rho_{\mathrm{ice}}}{\rho_{vs}^{I}}\sqrt{\frac{2\pi m}{k_B T}} ,
\label{eq:gt}
\end{equation}
with $\chi$ the curvature (positive for a convex grain). Here $\beta$ stands
for $\beta_{\mathrm{sub}} + \Delta\beta$; the thin-interface term is 2\% of
$\beta_{\mathrm{sub}}$ at every temperature in the aggregates
(Section~\ref{sec:neglected}), so it rescales $\beta$ by a constant.

The temperature-dependent quantities are $\rho_{vs}^{I}(T)$, $D_v(T)$,
$d_0(T) \propto 1/T$ and the explicit $\sqrt T$.

\section{Scales}

\begin{center}\begin{tabular}{lll}\toprule
Quantity & Scale & Reason\\\midrule
length, $\mathbf x = \bar R\,\mathbf x'$ & $\bar R$ & mean grain radius\\
curvature, $\chi = \chi'/\bar R$ & $1/\bar R$ & follows from the length\\
supersaturation, $\dfrac{\rho_v - \rho_{vs}^{I}}{\rho_{vs}^{I}} = \dfrac{d_0}{\bar R}\,u$ & $d_0/\bar R$ & what curvature $1/\bar R$ produces\\[1.2ex]
velocity, $v_n = V\,v'$ & $V = \dfrac{d_0}{\bar R\,\beta}$ & makes the kinetic term of order one\\[1.2ex]
time, $t = \tau\,t'$ & $\tau = \dfrac{\bar R}{V} = \dfrac{\bar R^{2}\beta}{d_0}$ & time to move one grain radius\\\bottomrule
\end{tabular}\end{center}

So $u$ is the supersaturation measured in units of the equilibrium
supersaturation of a mean grain, and $\rho_v = \rho_{vs}^{I}\,(1 + (d_0/\bar R)\,u)$.

\section{Substitution}

\emph{Gibbs--Thomson}, Eq.~\eqref{eq:gt}:
\[
\frac{d_0}{\bar R}\,u = d_0\,\frac{\chi'}{\bar R} + \beta\,\frac{d_0}{\bar R\,\beta}\,v'
\quad\Longrightarrow\quad
\boxed{\,u = \chi' + v'\,}
\]
No parameter is left.

\emph{Mass balance}, Eq.~\eqref{eq:stefan}. The left side is
$\rho_{\mathrm{ice}} V v'$. On the right, $\partial\rho_v/\partial n =
\rho_{vs}^{I}\,(d_0/\bar R)\,(1/\bar R)\,\partial u/\partial n'$. Therefore
\[
v' = \frac{D_v\,\rho_{vs}^{I}\,d_0}{\rho_{\mathrm{ice}}\,V\,\bar R^{2}}\,\frac{\partial u}{\partial n'}
   = \frac{D_v\,\rho_{vs}^{I}\,\beta}{\rho_{\mathrm{ice}}\,\bar R}\,\frac{\partial u}{\partial n'} ,
\]
using $V = d_0/(\bar R\beta)$. Define
\begin{equation}
L^{*} \equiv \frac{D_v\,\beta\,\rho_{vs}^{I}}{\rho_{\mathrm{ice}}}
       = \frac{D_v(T)}{\alpha_c}\sqrt{\frac{2\pi m}{k_B T}} ,
\qquad
\Lambda \equiv \frac{L^{*}}{\bar R} ,
\label{eq:Lstar}
\end{equation}
where the second form of $L^{*}$ follows from the definition of $\beta$ in
Eq.~\eqref{eq:gt}: \textbf{the saturation vapor density cancels.} Then
\[
\boxed{\,v' = \Lambda\,\frac{\partial u}{\partial n'}\,}
\]

\emph{Diffusion}, Eq.~\eqref{eq:diff}. Both sides carry the factor
$\rho_{vs}^{I} d_0/\bar R$:
\[
\frac{1}{\tau}\frac{\partial u}{\partial t'} = \frac{D_v}{\bar R^{2}}\nabla'^{2}u
\quad\Longrightarrow\quad
P\,\frac{\partial u}{\partial t'} = \nabla'^{2}u ,
\qquad
P \equiv \frac{\bar R^{2}}{D_v\,\tau} = \frac{d_0}{D_v\,\beta} .
\]
$P$ is the time for vapor to diffuse across a grain divided by the time for
the grain to change. It is $10^{-11}$ (Section~\ref{sec:neglected}), so the
vapor field is always in equilibrium with the current geometry:
\[
\boxed{\,\nabla'^{2}u = 0\,}
\]

\section{The dimensionless problem}

\begin{equation}
\nabla'^{2}u = 0 \ \text{ in the pore};\qquad
u = \chi' + v', \quad v' = \Lambda\,\frac{\partial u}{\partial n'} \ \text{ on the ice surface},
\end{equation}
with periodic boundaries. Eliminating $v'$ gives a single boundary condition,
$u - \Lambda\,\partial u/\partial n' = \chi'$.

Given a starting geometry, the evolution in $t' = t/\tau$ depends on
$\Lambda$ and on nothing else. Two limits make its meaning clear:
\begin{itemize}
\item $\Lambda \to \infty$: $u$ is uniform, attachment is the slow step,
  and each surface moves at $v' = \bar u - \chi'$ with $\bar u$ set by
  conservation of ice. The evolution in $t/\tau$ is then identical at
  every temperature.
\item $\Lambda \to 0$: $u = \chi'$ on the surface and diffusion is the slow
  step.
\end{itemize}
The aggregates have $\Lambda = 2.5$ to 2.9: attachment is the larger
resistance over distances shorter than $L^{*} \approx 135~\mu$m, which
covers transport between neighboring surfaces, but diffusion is not
negligible.

\section{Where temperature enters}

\begin{center}\begin{tabular}{lrrrrr}\toprule
 & $-40\,^{\circ}$C & $-30\,^{\circ}$C & $-20\,^{\circ}$C & $-10\,^{\circ}$C & $-5\,^{\circ}$C\\\midrule
$\rho_{vs}^{I}$ [$10^{-4}$ kg\,m$^{-3}$] & 1.055 & 3.123 & 8.487 & 21.39 & 33.11\\
$\tau$ [days] & 1713 & 591 & 222 & 89.8 & 58.5\\
$\tau/\tau(-5\,^{\circ}\mathrm{C})$ & 29.3 & 10.1 & 3.79 & 1.53 & 1\\
$L^{*}$ [$\mu$m] & 122.5 & 129.4 & 136.5 & 143.6 & 147.1\\
$\Lambda = L^{*}/\bar R$ & 2.45 & 2.59 & 2.73 & 2.87 & 2.94\\\bottomrule
\end{tabular}\end{center}

\begin{itemize}
\item \textbf{The time scale} carries a factor of 29.3 across the range:
  $\tau \propto \beta/d_0 \propto \sqrt T / \rho_{vs}^{I}(T)$ (since
  $d_0 \propto 1/T$ and $\beta \propto 1/(\rho_{vs}^{I}\sqrt T)$). The
  factor is $31.4$ from $\rho_{vs}^{I}$, times $1.072$ from $\sqrt T$,
  divided by $1.150$ from $d_0$.
\item \textbf{The one remaining group} changes by 20\%:
  $L^{*} \propto D_v/\sqrt T \propto T^{1.31}$.
\end{itemize}

So both temperature-dependent transport properties are accounted for. The
saturation vapor density sets the clock. The diffusivity survives only in
$\Lambda$, weakly, and its effect is further reduced because attachment is
the larger resistance. This is why curves at different temperatures
superpose when plotted against $t/\tau$, closely but not exactly: the
simulated rates follow $\tau_{\mathrm{sub}}$ to within 4\%. A 20\% change in
$\Lambda$ is the right size to produce that; it has not been tested as the
cause.

\textbf{Relation to $\tau_{\mathrm{sub}}$.} The manuscript's clock is
$\tau_{\mathrm{sub}} = \varepsilon^{2}(\beta_{\mathrm{sub}} +
\Delta\beta)/d_0 = (\varepsilon/\bar R)^{2}\,\tau$ with $\beta$ including
$\Delta\beta$. With $\varepsilon = \bar R/50$ fixed, the two differ by the
constant 2500, so ratios between temperatures are the same. Check at
$-20\,^{\circ}$C: $(1/50)^{2} \times 221.9~\mathrm{d} \times 1.020 =
7822$~s, the value in Table~2.

\section{What was neglected, with sizes}
\label{sec:neglected}

\begin{center}\begin{tabular}{p{0.40\linewidth}p{0.22\linewidth}p{0.28\linewidth}}\toprule
Neglected & Group & Size, $-40$ to $-5\,^{\circ}$C\\\midrule
Vapor storage (time derivative in Eq.~\ref{eq:diff}) & $P/\xi_v = d_0/(\xi_v D_v\beta)$ & $10^{-9}$ to $2\times10^{-8}$\\
Latent heat changing the local saturation & $G_T = \dfrac{\rho_{\mathrm{ice}} L_{\mathrm{sub}}^{2} m \bar R}{k_B T^{2} k_{\mathrm{ice}}\beta}$ & $10^{-4}$ to $2.5\times10^{-3}$\\[1.5ex]
$\rho_v$ against $\rho_{\mathrm{ice}}$ in Eq.~\eqref{eq:stefan} & $\rho_{vs}^{I}/\rho_{\mathrm{ice}}$ & $10^{-7}$ to $4\times10^{-6}$\\
Thin-interface term, as a change with $T$ & $\Delta\beta/\beta_{\mathrm{sub}}$ & 0.020 at every $T$\\\bottomrule
\end{tabular}\end{center}

\begin{itemize}
\item \emph{Storage.} The solver multiplies the diffusion and source terms
  of the vapor equation by $\xi_v = 10^{-3}$ together, which slows the
  vapor's relaxation by $10^{3}$ and leaves the balance~\eqref{eq:stefan}
  unchanged. The storage group becomes $P/\xi_v$, still negligible.
\item \emph{Latent heat.} Growth at velocity $V$ releases heat that changes
  the surface temperature by about $\rho_{\mathrm{ice}} L_{\mathrm{sub}} V
  \bar R / k_{\mathrm{ice}} \sim 10^{-7}$~K, which shifts the local
  saturation by $G_T$ in units of $u$. This group does depend strongly on
  temperature (it contains $\rho_{vs}^{I}$), but it is at most
  $2.5\times10^{-3}$. It is the only place the heat equation enters.
\item \emph{Thin-interface term.} $\Delta\beta$ and $\beta_{\mathrm{sub}}$
  are both proportional to $\rho_{\mathrm{ice}}/\rho_{vs}^{I}$, so their
  ratio is nearly independent of temperature and $\beta$ is rescaled by a
  constant 1.020.
\end{itemize}

\section{What this does and does not show}
\begin{itemize}
\item It shows that the collapse in $t/\tau_{\mathrm{sub}}$ follows from
  the structure of the model, for any geometry, as long as $\alpha_c$ does
  not depend on temperature. It is a property of the model.
\item It identifies what would break the collapse: an $\alpha_c(T)$ (it
  would enter $\Lambda$ as well as $\tau$); a second mechanism with its own
  temperature dependence; a temperature gradient, which makes
  $\rho_{vs}^{I}$ vary in space; or a much larger $\alpha_c$ or grain size,
  which lowers $\Lambda$ and makes the diffusivity matter.
\item It is a sharp-interface argument. The simulations add the check that
  the diffuse-interface model behaves this way: $k_{\mathrm{eff}}$ against
  SSA agrees across temperature to 0.27\%, rates to 4\%.
\end{itemize}

\end{document}
